% !TEX program = pdflatex
% AURION X 3.0 / APEX ENGINE — Documentación institucional (diagramas TikZ avanzados)
% Compilar: pdflatex AURION_X_APEX_ARCHITECTURE.tex  (×2)  |  Overleaf: pdfLaTeX
\documentclass[11pt,a4paper]{article}

\usepackage[utf8]{inputenc}
\usepackage[T1]{fontenc}
\usepackage[spanish]{babel}
\usepackage{geometry}
\geometry{left=2cm,right=2cm,top=2.2cm,bottom=2.2cm}

\usepackage{helvet}
\renewcommand{\familydefault}{\sfdefault}

\usepackage{graphicx}
\usepackage{booktabs}
\usepackage{enumitem}
\usepackage{hyperref}
\usepackage{xcolor}
\usepackage{tcolorbox}
\tcbuselibrary{skins,breakable}

\usepackage{pgfplots}
\pgfplotsset{compat=1.18}

\usepackage{tikz}
\usetikzlibrary{
  arrows.meta,
  positioning,
  shapes.geometric,
  shapes.symbols,
  shapes.multipart,
  fit,
  backgrounds,
  calc,
  decorations.pathreplacing,
  patterns.meta,
  fadings
}

%──────────────────────────────────────────────────────────────────────────────
% Paleta corporativa AURION
%──────────────────────────────────────────────────────────────────────────────
\definecolor{AXNavy}{RGB}{12,35,75}
\definecolor{AXNavyLight}{RGB}{32,72,140}
\definecolor{AXGold}{RGB}{212,168,55}
\definecolor{AXGoldDark}{RGB}{165,125,30}
\definecolor{AXSlate}{RGB}{72,88,108}
\definecolor{AXIce}{RGB}{235,242,252}
\definecolor{AXMint}{RGB}{210,245,232}
\definecolor{AXSand}{RGB}{255,248,232}
\definecolor{AXLavender}{RGB}{238,228,255}
\definecolor{AXRose}{RGB}{255,228,232}
\definecolor{AXTeal}{RGB}{215,248,242}
\definecolor{AXGraphite}{RGB}{45,52,64}

\colorlet{AXLayer1}{AXIce}
\colorlet{AXLayer2}{AXMint}
\colorlet{AXLayer3}{AXSand}
\colorlet{AXLayer4}{AXLavender}
\colorlet{AXLayer5}{AXRose}
\colorlet{AXLayer6}{AXTeal}
\colorlet{AXLayer7}{red!8!white}
\colorlet{AXLayer8}{blue!6!white}
\colorlet{AXLayer9}{AXGraphite!12}

\hypersetup{
  colorlinks=true,
  linkcolor=AXNavyLight,
  urlcolor=AXNavyLight
}

\newtcolorbox{axcallout}[1][]{
  enhanced, breakable,
  colback=AXIce, colframe=AXNavyLight,
  fonttitle=\bfseries\sffamily,
  attach boxed title to top left={yshift=-2mm,xshift=4mm},
  boxed title style={colback=AXNavy, colframe=AXNavy, sharp corners},
  #1
}

%──────────────────────────────────────────────────────────────────────────────
% Estilos TikZ profesionales
%──────────────────────────────────────────────────────────────────────────────
\tikzset{
  axarr/.style={
    -{Stealth[length=3.2mm,width=2.4mm]},
    line width=1.1pt,
    draw=AXNavyLight!85,
    rounded corners=2pt
  },
  axarrgold/.style={
    -{Stealth[length=3mm,width=2.2mm]},
    line width=1pt,
    dashed,
    draw=AXGoldDark!90
  },
  axlbl/.style={
    font=\footnotesize\sffamily\bfseries,
    text=AXSlate,
    fill=white,
    inner sep=2pt,
    rounded corners=1pt,
    opacity=0.95
  },
  axlblw/.style={
    font=\scriptsize\sffamily,
    text=white,
    fill=AXNavy!85,
    inner sep=2.5pt,
    rounded corners=2pt
  },
  axpipe/.style={
    circle, fill=AXNavyLight, inner sep=0pt,
    minimum size=5pt
  }
}

% Caja de capa con barra lateral (estilo enterprise)
% Uso: \LayerBoxAt{coord}{ancho}{alto}{id}{color}{titulo}{subtitulo}{numero}
\newcommand{\LayerBoxAt}[8]{%
  \node[
    anchor=west, inner sep=0pt,
    minimum width=#2, minimum height=#3
  ] (#4-outer) at #1 {};
  \fill[#5!88!white] (#4-outer.south west) rectangle (#4-outer.north east);
  \fill[#5] (#4-outer.south west) rectangle ([xshift=5.5mm]#4-outer.south west |- #4-outer.north west);
  \draw[AXNavy!50, line width=0.6pt] (#4-outer.south west) rectangle (#4-outer.north east);
  \draw[AXNavy!70, line width=1.2pt] ([xshift=5.5mm]#4-outer.south west) -- ([xshift=5.5mm]#4-outer.north west);
  \node[anchor=west, font=\bfseries\footnotesize\sffamily, text=white]
    at ([xshift=1.3mm]#4-outer.west |- #4-outer.center) {#8};
  \node[
    anchor=west, align=left, text width=#2-14mm,
    font=\small\sffamily, text=AXGraphite
  ] at ([xshift=7mm]#4-outer.west |- #4-outer.center) {%
    \textbf{#6}\\[-1pt]
    {\footnotesize #7}%
  };
}

% Portada
\newcommand{\makeaxtitle}{%
  \begin{titlepage}
    \begin{tikzpicture}[remember picture,overlay]
      \fill[AXNavy] (current page.south west) rectangle (current page.north east);
      \fill[AXNavyLight!40] ([yshift=4cm]current page.south west) rectangle ([yshift=8cm]current page.south east);
      \fill[AXGold] ([yshift=7.85cm]current page.south west) rectangle ([yshift=8cm]current page.south east);
      \foreach \i in {0,...,18} {
        \draw[white, opacity=0.04, line width=0.5pt]
          ([xshift=\i cm]current page.south west) -- ++(45:18cm);
      }
      \node[anchor=west, text=white] at ([xshift=2.2cm,yshift=10cm]current page.south west) {%
        {\fontsize{42}{46}\selectfont\bfseries AURION\_X 3.0}\\[4pt]
        {\fontsize{22}{26}\selectfont APEX ENGINE}\\[14pt]
        {\large Arquitectura del sistema cuantitativo institucional}\\[6pt]
        {\normalsize Documentación ejecutiva y técnica · Versión 3.0}
      };
      \node[
        draw=AXGold, line width=1.5pt, rounded corners=6pt,
        fill=AXNavy!60, text=white, inner sep=10pt,
        font=\footnotesize\sffamily
      ] at ([xshift=2.2cm,yshift=3.2cm]current page.south west) {%
        XAUUSD · USTEC · BTCUSD \quad|\quad PostgreSQL · MT5 · ML Ensemble
      };
    \end{tikzpicture}
  \end{titlepage}
}

\pagestyle{plain}
\setcounter{secnumdepth}{2}

\begin{document}
\hypersetup{pageanchor=false}
\makeaxtitle
\hypersetup{pageanchor=true}
\tableofcontents
\newpage

%==============================================================================
\section{Visión general del proyecto}
%==============================================================================

\begin{axcallout}{title=Resumen ejecutivo}
\textbf{AURION\_X 3.0} es un motor cuantitativo de trading algorítmico que opera en \textbf{MetaTrader 5}, aprende de resultados históricos en \textbf{PostgreSQL} y entra al mercado solo cuando existe \textbf{ventaja estadística + contexto institucional} verificable por ocho capas de decisión.
\end{axcallout}

\subsection{Stack de capacidades}
\begin{center}
\begin{tikzpicture}[
  cap/.style={
    rectangle, rounded corners=5pt,
    minimum width=3.1cm, minimum height=1.05cm,
    align=center, font=\footnotesize\sffamily\bfseries,
    draw=#1!60, fill=#1!15, text=AXGraphite,
    opacity=0.98
  }
]
  \node[cap=AXNavyLight] (a) {Trading\\algorítmico};
  \node[cap=AXLavender, right=3mm of a] (b) {ML\\adaptativo};
  \node[cap=AXSand, right=3mm of b] (c) {Análisis\\institucional};
  \node[cap=AXMint, right=3mm of c] (d) {Micro-\\estructura};
  \node[cap=AXRose, below=4mm of b] (e) {Riesgo\\prop firm};
  \node[cap=AXTeal, right=3mm of e] (f) {Auto-\\aprendizaje};
  \draw[axarr] (a) -- (b) -- (c) -- (d);
  \draw[axarr] (b) -- (e) -- (f);
\end{tikzpicture}
\end{center}

\subsection{Objetivos de construcción}
Detectar oportunidades de alta probabilidad · minimizar drawdown · adaptación automática · operación institucional · estabilidad y consistencia en cuentas funded.

%==============================================================================
\section{Objetivo principal y filosofía}
%==============================================================================

\begin{center}
\begin{tikzpicture}
  \node[
    rectangle, rounded corners=8pt,
    minimum width=0.92\textwidth, minimum height=1.4cm,
    draw=AXGold, line width=1.5pt,
    fill=AXSand!50
  ] (q) {};
  \node[text=AXNavy, font=\large\sffamily\itshape] at (q.center) {%
    ``Operar únicamente cuando existe una ventaja estadística y contextual clara.''
  };
\end{tikzpicture}
\end{center}

\begin{figure}[htbp]
\centering
\begin{tikzpicture}[node distance=1.4cm]
  \node[
    rectangle, rounded corners=6pt, minimum width=4.8cm, minimum height=2.2cm,
    draw=red!50!black, fill=red!6  ] (no) {};
  \node[text=AXGraphite, align=center] at (no.center) {%
    \textbf{\color{red!70!black}ENFOQUE RETAIL}\\[4pt]
    \footnotesize Rupturas ciegas\\
    \footnotesize Perseguir velas\\
    \footnotesize Operar ruido / chop
  };
  \node[
    rectangle, rounded corners=6pt, minimum width=4.8cm, minimum height=2.2cm,
    draw=AXNavyLight!70, fill=AXMint!40,
    right=2.2cm of no
  ] (si) {};
  \node[text=AXGraphite, align=center] at (si.center) {%
    \textbf{\color{AXNavy}ENFOQUE INSTITUCIONAL}\\[4pt]
    \footnotesize Liquidez y sweeps\\
    \footnotesize Absorción y manipulación\\
    \footnotesize Desplazamiento real
  };
  \node[axlbl, above=6mm of no] {Evita};
  \node[axlbl, above=6mm of si] {Prioriza};
  \draw[axarrgold, line width=1.4pt] (no) -- node[axlbl, above, sloped] {evolución} (si);
\end{tikzpicture}
\caption{Filosofía operativa del motor AURION.}
\end{figure}

%==============================================================================
\section{Activos y despliegue multi-símbolo}
%==============================================================================

\begin{center}
\begin{tikzpicture}[
  asset/.style={
    circle, minimum size=2.1cm, align=center,
    font=\footnotesize\sffamily\bfseries, text=white
  }
]
  \node[asset, fill=AXGold!90!black] (xau) {XAUUSD\\{\tiny Oro}};
  \node[asset, fill=AXNavyLight, right=1.8cm of xau] (ustec) {USTEC\\{\tiny Índice}};
  \node[asset, fill=AXGraphite, right=1.8cm of ustec] (btc) {BTCUSD\\{\tiny Cripto}};
  \node[
    rectangle, rounded corners=4pt, draw=AXNavyLight, dashed,
    fit=(xau)(btc), inner sep=14pt,
    fill=AXIce!30, label={[axlbl]above:{\textbf{Un proceso · Un orchestrator · Tres modelos PKL}}}
  ] {};
  \draw[axarr] (xau) -- node[axlbl, above] {pesos ML} (ustec);
  \draw[axarr] (ustec) -- node[axlbl, above] {memoria} (btc);
\end{tikzpicture}
\end{center}

Cada activo: \texttt{model\_SYMBOL.pkl} · \texttt{trade\_memory\_SYMBOL.json} · pesos en \texttt{asset\_profiles}.

%==============================================================================
\section{Arquitectura en 9 capas}
%==============================================================================

\begin{figure}[htbp]
\centering
\begin{tikzpicture}[
  layerbox/.style={
    rectangle, rounded corners=4pt, anchor=west,
    minimum width=11.2cm, minimum height=1.05cm,
    align=left, text width=10.2cm, inner sep=6pt,
    font=\small\sffamily, draw=AXNavy!55, line width=0.8pt
  },
  badge/.style={
    circle, minimum size=6.5mm, fill=AXNavy, text=white,
    font=\scriptsize\bfseries\sffamily, inner sep=0pt
  }
]
  \node[layerbox, fill=AXLayer1] (L1) {%
    \tikz[baseline=-0.5ex]{\node[badge] at (0,0) {1};}\hspace{0.35cm}%
    \textbf{Capa 1 — Market Data}\hfill{\footnotesize MT5 · TickHandler · MTF}};
  \node[layerbox, fill=AXLayer2, below=3.5mm of L1] (L2) {%
    \tikz[baseline=-0.5ex]{\node[badge] at (0,0) {2};}\hspace{0.35cm}%
    \textbf{Capa 2 — Market Analysis}\hfill{\footnotesize Regime · Microstructure}};
  \node[layerbox, fill=AXLayer3, below=3.5mm of L2] (L3) {%
    \tikz[baseline=-0.5ex]{\node[badge] at (0,0) {3};}\hspace{0.35cm}%
    \textbf{Capa 3 — Institutional}\hfill{\footnotesize Sweeps · Absorción}};
  \node[layerbox, fill=AXLayer4, below=3.5mm of L3] (L4) {%
    \tikz[baseline=-0.5ex]{\node[badge] at (0,0) {4};}\hspace{0.35cm}%
    \textbf{Capa 4 — Machine Learning}\hfill{\footnotesize Ensemble · Quant}};
  \node[layerbox, fill=AXLayer5, below=3.5mm of L4] (L5) {%
    \tikz[baseline=-0.5ex]{\node[badge] at (0,0) {5};}\hspace{0.35cm}%
    \textbf{Capa 5 — Decision (APEX)}\hfill{\footnotesize Hierarchy · Hyper}};
  \node[layerbox, fill=AXLayer6, below=3.5mm of L5] (L6) {%
    \tikz[baseline=-0.5ex]{\node[badge] at (0,0) {6};}\hspace{0.35cm}%
    \textbf{Capa 6 — Execution}\hfill{\footnotesize Exit AI · Orders}};
  \node[layerbox, fill=AXLayer7, below=3.5mm of L6] (L7) {%
    \tikz[baseline=-0.5ex]{\node[badge] at (0,0) {7};}\hspace{0.35cm}%
    \textbf{Capa 7 — Risk}\hfill{\footnotesize FundingPips · MetaRisk}};
  \node[layerbox, fill=AXLayer8, below=3.5mm of L7] (L8) {%
    \tikz[baseline=-0.5ex]{\node[badge] at (0,0) {8};}\hspace{0.35cm}%
    \textbf{Capa 8 — Self-Learning}\hfill{\footnotesize Retrain · Evolution}};
  \node[layerbox, fill=AXLayer9, below=3.5mm of L8] (L9) {%
    \tikz[baseline=-0.5ex]{\node[badge] at (0,0) {9};}\hspace{0.35cm}%
    \textbf{Capa 9 — Database}\hfill{\footnotesize PostgreSQL · Dashboard}};

  \draw[axarr, line width=1.2pt] ([xshift=5.8cm]L1.north) -- ([xshift=5.8cm]L9.south);
  \node[axlblw, rotate=-90] at ([xshift=6.15cm]$(L1.north)!0.5!(L9.south)$) {Flujo};

  \foreach \a/\b in {L1/L2,L2/L3,L3/L4,L4/L5,L5/L6,L6/L7,L7/L8,L8/L9} {
    \draw[axarr] (\a.north) -- (\b.south);
  }
  \begin{scope}[on background layer]
    \node[
      draw=AXGold, line width=1.2pt, dashed, rounded corners=8pt,
      fit=(L3)(L5), inner sep=9pt, fill=AXGold!10
    ] {};
    \node[axlblw, anchor=north west] at ([xshift=2mm,yshift=-1mm]L3.north west) {CORE};
  \end{scope}
\end{tikzpicture}
\caption{Arquitectura modular en nueve capas (optimizada para impresión PDF).}
\end{figure}

%==============================================================================
\section{Hub central: Orchestrator}
%==============================================================================

\begin{figure}[htbp]
\centering
\begin{tikzpicture}[
  hub/.style={
    rectangle, rounded corners=8pt,
    minimum width=3.4cm, minimum height=1.1cm,
    align=center, font=\footnotesize\sffamily,
    draw=AXNavyLight!70, fill=white,
  },
  core/.style={
    circle, minimum size=2.8cm, align=center,
    font=\small\sffamily\bfseries, text=white,
    fill=AXNavy  }
]
  \node[core] (orch) {ORCHESTRATOR\\[-2pt]{\footnotesize core/main}};
  \node[hub, fill=AXIce, above left=1.1cm and 0.9cm of orch] (mt5) {\textbf{MT5}\\Connector};
  \node[hub, fill=AXMint, above right=1.1cm and 0.9cm of orch] (strat) {\textbf{Strategy}\\Scalper · Intraday};
  \node[hub, fill=AXLavender, left=1.5cm of orch] (apex) {\textbf{ApexEngine}\\Institutional};
  \node[hub, fill=AXRose, right=1.5cm of orch] (risk) {\textbf{Risk}\\Manager};
  \node[hub, fill=AXSand, below left=1.1cm and 0.9cm of orch] (pg) {\textbf{PostgreSQL}\\TradeRepo};
  \node[hub, fill=AXTeal, below right=1.1cm and 0.9cm of orch] (exec) {\textbf{Execution}\\Orders};

  \foreach \n in {mt5,strat,apex,risk,pg,exec} {
    \draw[axarr, line width=1.2pt] (\n) -- (orch);
  }
  \draw[axarrgold] (exec.east) -- ++(0.6,0) |- (mt5.east)
    node[pos=0.25, axlbl, below] {ciclo tick};

  \begin{scope}[on background layer]
    \node[
      draw=AXNavy!30, dashed, rounded corners=12pt,
      fit=(mt5)(strat)(apex)(risk)(pg)(exec)(orch),
      inner sep=16pt, fill=AXIce!20
    ] {};
  \end{scope}
\end{tikzpicture}
\caption{Topología hub-and-spoke: un único proceso Python, multi-símbolo.}
\end{figure}

%==============================================================================
\section{Jerarquía de decisión Apex (8 capas)}
%==============================================================================

\begin{figure}[htbp]
\centering
\begin{tikzpicture}[
  funnel/.style={
    trapezium, trapezium left angle=78, trapezium right angle=78,
    draw=AXNavy!60, line width=0.9pt,
    minimum height=0.95cm, text width=9.2cm, align=center,
    font=\small\sffamily,
    #1
  }
]
  \node[funnel, fill=red!12] (f1) {%
    \textbf{\textcolor{red!65!black}{1 · RIESGO GLOBAL}}\\
    {\footnotesize Drawdown · Hard stop · MetaRisk DEFENSIVE/AGGRESSIVE}};
  \node[funnel, fill=orange!15, below=3mm of f1] (f2) {%
    \textbf{2 · CONTEXTO HIPER} — HyperContextEngine (fusión multi-señal)};
  \node[funnel, fill=yellow!18, below=3mm of f2] (f3) {%
    \textbf{3 · RÉGIMEN / ESTADO LATENTE} — mercado peligroso · chop};
  \node[funnel, fill=AXSand!80, below=3mm of f3] (f4) {%
    \textbf{4 · MICROESTRUCTURA} — calidad · manipulación · velocidad};
  \node[funnel, fill=AXLavender!70, below=3mm of f4] (f5) {%
    \textbf{5 · MACHINE LEARNING} — umbral probabilístico · ensemble v2};
  \node[funnel, fill=AXRose!50, below=3mm of f5] (f6) {%
    \textbf{6 · SETUP RANKING} — grade A+ … X · memoria histórica};
  \node[funnel, fill=AXTeal!60, below=3mm of f6] (f7) {%
    \textbf{7 · CALIDAD DE EJECUCIÓN} — spread · ExecIntelligence};
  \node[funnel, fill=AXMint!70, below=3mm of f7] (f8) {%
    \textbf{8 · ENTRADA FINAL} — anti-overtrade · spacing · burst};

  \node[
    diamond, aspect=2.2, draw=AXGold, line width=1.8pt,
    fill=AXGold!25, minimum width=5.5cm, align=center,
    below=8mm of f8, font=\small\sffamily\bfseries  ] (out) {ALLOW TRADE\\[-2pt]{\footnotesize o BLOCK + razón de capa}};

  \foreach \a/\b in {f1/f2,f2/f3,f3/f4,f4/f5,f5/f6,f6/f7,f7/f8,f8/out} {
    \draw[axarr] (\a) -- (\b);
  }
\end{tikzpicture}
\caption{Embudo de decisión institucional. Si la capa $n$ falla, las capas $n{+}1\ldots 8$ ni siquiera se evalúan.}
\end{figure}

%==============================================================================
\section{Flujo operativo completo (12 pasos)}
%==============================================================================

\begin{figure}[htbp]
\centering
\resizebox{\textwidth}{!}{%
\begin{tikzpicture}[
  step/.style={
    rectangle, rounded corners=4pt,
    minimum width=2.15cm, minimum height=1.35cm,
    align=center, font=\scriptsize\sffamily,
    draw=#1!70, fill=#1!12, text=AXGraphite
  }
]
  % Swimlane DATA
  \node[step=AXIce, anchor=west] (s1) at (0,0) {\textbf{01}\\MT5\\Ticks};
  \node[step=AXIce, right=2mm of s1] (s2) {\textbf{02}\\Market\\Layer};
  \node[step=AXSand, right=2mm of s2] (s3) {\textbf{03}\\Instit.\\Analysis};

  % Swimlane INTELLIGENCE
  \node[step=AXLavender, right=2mm of s3] (s4) {\textbf{04}\\ML\\Predict};
  \node[step=AXRose, right=2mm of s4] (s5) {\textbf{05}\\Apex\\Fusion};
  \node[step=AXRose, right=2mm of s5] (s6) {\textbf{06}\\Hierarchy\\8 capas};

  % Swimlane CONTROL
  \node[step=AXRose, right=2mm of s6] (s7) {\textbf{07}\\Risk\\Gate};

  % Swimlane EXEC + LEARN
  \node[step=AXTeal, below=1.1cm of s1] (s8) {\textbf{08}\\Execute\\Order};
  \node[step=AXTeal, right=2mm of s8] (s9) {\textbf{09}\\Exit\\AI};
  \node[step=AXGraphite, right=2mm of s9] (s10) {\textbf{10}\\PostgreSQL\\Persist};
  \node[step=AXNavyLight, right=2mm of s10] (s11) {\textbf{11}\\Self-\\Learning};
  \node[step=AXGoldDark, right=2mm of s11] (s12) {\textbf{12}\\Auto\\Retrain};

  \foreach \x/\y in {s1/s2,s2/s3,s3/s4,s4/s5,s5/s6,s6/s7} {
    \draw[axarr] (\x) -- (\y);
  }
  \draw[axarr] (s7.south) ++(0,-0.15) -| (s8.north);
  \foreach \x/\y in {s8/s9,s9/s10,s10/s11,s11/s12} {
    \draw[axarr] (\x) -- (\y);
  }
  \draw[axarrgold] (s12.north) -- ++(0,0.55) -| node[axlbl, near start, above] {feedback modelos} (s4.north);

  % Swimlane labels
  \node[rotate=90, axlbl, anchor=south] at (-0.85,0.55) {DATOS};
  \node[rotate=90, axlbl, anchor=south] at (-0.85,-1.35) {EJECUCIÓN \& APRENDIZAJE};

  \begin{scope}[on background layer]
    \draw[AXNavy!15, fill=AXIce!25, rounded corners=6pt]
      (-1.1,-0.75) rectangle (16.2,1.15);
    \draw[AXNavy!15, fill=AXMint!15, rounded corners=6pt]
      (-1.1,-2.35) rectangle (16.2,-0.55);
  \end{scope}
\end{tikzpicture}%
}
\caption{Pipeline end-to-end con swimlanes y retroalimentación de retrain.}
\end{figure}

%==============================================================================
\section{Pipeline Machine Learning}
%==============================================================================

\begin{figure}[htbp]
\centering
\begin{tikzpicture}[
  proc/.style={
    rectangle, rounded corners=5pt,
    minimum width=2.6cm, minimum height=1.2cm,
    align=center, font=\footnotesize\sffamily,
    draw=AXNavyLight, fill=white,
    double=AXNavyLight, double distance=1.5pt
  }
]
  \node[proc, fill=AXIce!50] (raw) {\textbf{Raw Data}\\MT5 + PG trades};
  \node[proc, fill=AXSand!40, right=1.1cm of raw] (fe) {\textbf{Features v2}\\25+ variables};
  \node[proc, fill=AXLavender!40, right=1.1cm of fe] (ens) {\textbf{Ensemble}\\RF·XGB·LGBM};
  \node[proc, fill=AXRose!30, right=1.1cm of ens] (inf) {\textbf{Inference}\\$P_{buy}$ · $P_{sell}$};

  \draw[axarr] (raw) -- node[axlbl, above] {build} (fe);
  \draw[axarr] (fe) -- node[axlbl, above] {train} (ens);
  \draw[axarr] (ens) -- node[axlbl, above] {predict} (inf);

  \node[
    rectangle, rounded corners=5pt, draw=AXGold!80!black, thick,
    fill=AXGold!12, minimum width=3.2cm, minimum height=2cm,
    below=1.3cm of ens, align=center, font=\footnotesize\sffamily
  ] (rt) {\textbf{Retrain Loop}\\[2pt]
    Walk-forward · OOS gate\\
    Rollback · Hot-reload};

  \draw[axarrgold] (inf.south) ++(0,-0.2) -| (rt.north);
  \draw[axarrgold] (rt.west) -| node[axlbl, below, pos=0.3] {deploy PKL} (ens.south);
  \draw[axarrgold] (raw.south) |- (rt.west);

  \node[axlblw, above=2mm of inf] {QuantDecision};
\end{tikzpicture}
\caption{Pipeline ML con doble borde, bucle de retrain y validación OOS.}
\end{figure}

\textbf{Modelos:} RandomForest · LogisticRegression · XGBoost · LightGBM · Voting Ensemble.

%==============================================================================
\section{Base de datos PostgreSQL}
%==============================================================================

\begin{figure}[htbp]
\centering
\begin{tikzpicture}[
  db/.style={
    rectangle, rounded corners=6pt, draw=AXNavy, line width=1.2pt,
    fill=AXNavyLight!15, minimum width=3cm, minimum height=1.4cm,
    align=center, font=\footnotesize\sffamily\bfseries, text=AXNavy
  },
  tab/.style={
    rectangle, draw=AXNavyLight!70, rounded corners=4pt,
    minimum width=3.4cm, minimum height=1.15cm, fill=white,
    align=center, font=\small\sffamily
  }
]
  \node[db] (pg) {PostgreSQL\\{\tiny Docker · aurionx}};

  \node[tab, above right=0.85cm and 1.1cm of pg] (pred) {%
    \textbf{ml\_predictions}\\{\scriptsize conf · probs · grade}};
  \node[tab, right=0.65cm of pred] (feat) {%
    \textbf{ml\_features}\\{\scriptsize vector v2 · sweeps}};
  \node[tab, below right=0.85cm and 1.1cm of pg] (reg) {%
    \textbf{ml\_regime}\\{\scriptsize vol · sesión · mode}};
  \node[tab, right=0.65cm of reg] (perf) {%
    \textbf{ml\_performance}\\{\scriptsize sharpe · expectancy}};

  \draw[axarr] (pg) -- (pred);
  \draw[axarr] (pg) -- (feat);
  \draw[axarr] (pg) -- (reg);
  \draw[axarr] (pg) -- (perf);
  \draw[axarr] (pred) -- node[axlbl, above] {FK} (feat);
  \draw[axarrgold] (perf) to[bend left=30] node[axlbl, right] {retrain} (pred);
\end{tikzpicture}
\caption{Modelo lógico de datos PostgreSQL y tablas ML.}
\end{figure}

%==============================================================================
\section{Auto-aprendizaje y ciclo de evolución}
%==============================================================================

\begin{figure}[htbp]
\centering
\begin{tikzpicture}[
  phase/.style={
    circle, minimum size=1.55cm, align=center,
    font=\footnotesize\sffamily\bfseries, text=white
  }
]
  \node[phase, fill=AXNavyLight] (p1) at (0,0) {1\\Gather};
  \node[phase, fill=AXNavyLight!85!white] (p2) at ({60}:2.8cm) {2\\Analyze};
  \node[phase, fill=AXGold!90!black] (p3) at ({120}:2.8cm) {3\\Patterns};
  \node[phase, fill=AXGold!70!black] (p4) at ({180}:2.8cm) {4\\Weights};
  \node[phase, fill=AXRose!70!black] (p5) at ({240}:2.8cm) {5\\Train};
  \node[phase, fill=AXRose!50!black] (p6) at ({300}:2.8cm) {6\\Validate};
  \node[phase, fill=AXMint!70!black] (p7) at ({0}:2.8cm) {7\\Deploy};

  \foreach \a/\b in {p1/p2,p2/p3,p3/p4,p4/p5,p5/p6,p6/p7,p7/p1} {
    \draw[axarr, line width=1.3pt] (\a) -- (\b);
  }
  \node[align=center, font=\small\sffamily\bfseries, text=AXNavy] at (0,0) {%
    AUTO\\EVOLVE
  };
  \node[axlbl, below=3.6cm of p4] {Cada 10 cierres · Diario 22:00 UTC · Degradación · Cambio régimen};
\end{tikzpicture}
\caption{Ciclo cerrado de auto-aprendizaje (diagrama radial).}
\end{figure}

%==============================================================================
\section{Gestión de riesgo y prop firm}
%==============================================================================

\begin{figure}[htbp]
\centering
\begin{tikzpicture}
  \node[
    regular polygon, regular polygon sides=6, minimum size=3.2cm,
    draw=AXNavy, line width=1.5pt, fill=AXRose!25,
    align=center, font=\footnotesize\sffamily  ] (shield) at (0,0) {%
    \textbf{RISK}\\[-2pt]\textbf{SHIELD}
  };

  \node[
    rectangle, rounded corners=4pt, minimum width=2.9cm, minimum height=0.95cm,
    draw=AXGold!80!black, fill=AXSand!50, align=center, font=\scriptsize\sffamily,
    above left=1.2cm and 0.3cm of shield
  ] (fp) {\textbf{FundingPips}\\DD diario · DD total};
  \node[rectangle, rounded corners=4pt, minimum width=2.9cm, minimum height=0.95cm,
    draw=AXNavyLight, fill=AXLavender!30, align=center, font=\scriptsize\sffamily,
    above right=1.2cm and 0.3cm of shield] (mr) {\textbf{MetaRisk}\\DEF · AGG · NEUT};
  \node[rectangle, rounded corners=4pt, minimum width=2.9cm, minimum height=0.95cm,
    draw=red!50, fill=red!8, align=center, font=\scriptsize\sffamily,
    below left=1.2cm and 0.3cm of shield] (ao) {\textbf{Anti-OT}\\spacing · burst};
  \node[rectangle, rounded corners=4pt, minimum width=2.9cm, minimum height=0.95cm,
    draw=AXMint!70!black, fill=AXMint!30, align=center, font=\scriptsize\sffamily,
    below right=1.2cm and 0.3cm of shield] (mc) {\textbf{Monte Carlo}\\lot factor};

  \foreach \n in {fp,mr,ao,mc} {
    \draw[axarr] (\n) -- (shield);
  }

  \node[
    rectangle, rounded corners=6pt, draw=AXNavy, line width=1.2pt,
    fill=AXGold!20, minimum width=4.5cm, minimum height=0.9cm,
    below=1.8cm of shield, font=\small\sffamily\bfseries  ] (gate) {GATE DE ENTRADA — approve / block};

  \draw[axarr, line width=1.4pt] (shield) -- (gate);
\end{tikzpicture}
\caption{Escudo de riesgo multicapa antes de \texttt{order\_send}.}
\end{figure}

%==============================================================================
\section{Diagrama de secuencia (tick $\rightarrow$ trade)}
%==============================================================================

\begin{figure}[htbp]
\centering
\begin{tikzpicture}[
  life/.style={rectangle, draw=AXNavyLight, fill=AXIce!40, minimum width=1.35cm,
    minimum height=0.65cm, font=\scriptsize\sffamily\bfseries},
  msg/.style={-{Stealth[length=2mm]}, line width=0.9pt, draw=AXNavyLight}
]
  \node[life] (mt5) at (0,0) {MT5};
  \node[life, right=1.5cm of mt5] (orch) {Orch.};
  \node[life, right=1.5cm of orch] (apex) {Apex};
  \node[life, right=1.5cm of apex] (risk) {Risk};
  \node[life, right=1.5cm of risk] (brok) {Broker};

  \foreach \n in {mt5,orch,apex,risk,brok} {
    \draw[AXNavyLight!40, dashed] (\n.south) -- ++(0,-5.2);
  }

  \draw[msg] ([yshift=-0.8cm]mt5.east) -- node[axlbl, above] {tick} ([yshift=-0.8cm]orch.west);
  \draw[msg] ([yshift=-1.5cm]orch.east) -- node[axlbl, above] {evaluate} ([yshift=-1.5cm]apex.west);
  \draw[msg] ([yshift=-2.2cm]apex.east) -- node[axlbl, above] {allow?} ([yshift=-2.2cm]risk.west);
  \draw[msg] ([yshift=-3.0cm]orch.east) -- node[axlbl, above] {order} ([yshift=-3.0cm]brok.west);
  \draw[msg, dashed] ([yshift=-3.8cm]brok.west) -- node[axlbl, below] {fill} ([yshift=-3.8cm]orch.east);
  \draw[msg, axarrgold] ([yshift=-4.6cm]orch.south) ++(-0.5,0) -- ++(-0.5,-0.5)
    node[axlbl, left] {persist PG} ++(1,0);

  \node[axlblw, below=5.6cm of brok] {Secuencia simplificada en vivo};
\end{tikzpicture}
\caption{Interacción temporal entre componentes en un ciclo de entrada.}
\end{figure}

%==============================================================================
\section{Microestructura, salidas y régimenes}
%==============================================================================

\textbf{Microestructura:} sweeps · fake breakouts · absorción · manipulación · liquidity grabs · stop hunts.

\textbf{TradeExitAI:} momentum loss · agotamiento · absorción contraria · compresión · partial TP.

\textbf{Régimenes:} trend · range · breakout · defensive · high volatility — transición vía \texttt{MarketRegimeSwitcher}.

%==============================================================================
\section{Dashboard y tecnologías}
%==============================================================================

Dashboard HTTP (\texttt{127.0.0.1:8765}): equity · winrate · setups · drawdown · confianza ML.

\begin{center}
\begin{tabular}{@{}cccc@{}}
\toprule
\textbf{Python} & \textbf{MT5} & \textbf{PostgreSQL} & \textbf{Sklearn / XGB / LGBM} \\
\midrule
Docker & Pandas & NumPy & SQLAlchemy \\
\bottomrule
\end{tabular}
\end{center}

%==============================================================================
\section{Comparativa y conclusión}
%==============================================================================

\begin{figure}[htbp]
\centering
\begin{tikzpicture}
  \begin{axis}[
    ybar,
    bar width=14pt,
    ymin=0, ymax=100,
    ytick={0,25,50,75,100},
    symbolic x coords={Contexto,ML,Memoria,Riesgo,Adapt.},
    xtick=data,
    width=11cm, height=5cm,
    axis lines*=left,
    tick label style={font=\footnotesize\sffamily},
    nodes near coords,
    nodes near coords align={vertical},
    every node near coord/.append style={font=\tiny\sffamily}
  ]
  \addplot[fill=red!25, draw=red!50] coordinates {
    (Contexto,18) (ML,22) (Memoria,12) (Riesgo,15) (Adapt.,10)
  };
  \addplot[fill=AXNavyLight!60, draw=AXNavy] coordinates {
    (Contexto,92) (ML,88) (Memoria,95) (Riesgo,90) (Adapt.,87)
  };
  \legend{Bot retail típico,AURION X 3.0}
  \end{axis}
\end{tikzpicture}
\caption{Madurez relativa por dimensión (índice cualitativo 0–100).}
\end{figure}

\begin{axcallout}{title=Conclusión}
\textbf{AURION\_X 3.0} evoluciona de un EA tradicional hacia un \textbf{sistema cuantitativo institucional auto-adaptativo}: aprende, protege capital, fusiona contexto profundo y mejora continuamente — arquitectura modular orientada a consistencia y supervivencia estadística en mercados reales y prop firms.
\end{axcallout}

\vfill
\begin{center}
\footnotesize\sffamily
\textcolor{AXNavy}{\textbf{AURION X 3.0 / APEX ENGINE}}\\
\texttt{docs/AURION\_X\_APEX\_ARCHITECTURE.tex} · Compilar con pdfLaTeX o Overleaf
\end{center}

\end{document}
